\section{Introduction}
\IEEEPARstart{P}{lanned} self-tests can deliberately move a controlled mechanism between configurations to discriminate a specified fault alternative from admissible healthy drift. This paper studies that design problem over a fixed observation horizon. The task calendar and fault onset are specified before data are observed; the objective is fixed-horizon minimax detection power, with optional preassigned multiple-test error allocation. The design is useful when a dominant healthy response reverses between two tasks while an additive fault response does not. A transition consumes observation time and permits the healthy parameter to change, so scheduling changes the uncertainty set as well as the sampling budget.

We determine the optimal information deficit over all allowed deterministic calendars. For linear fault growth with an unbounded healthy level, the deficit is asymptotic to a known constant times $H^{5/2}$. A schedule-uniform converse matches a symmetric half-window--full-window--half-window construction. For a fixed finite healthy interval, the deficit instead equals an explicit quadratic polynomial up to $O(1)$, and switching has only a bounded advantage over staying. The change of regime follows from the same physical-time drift constraint.

The system-design consequence is quantitative: transition time $\tau_m$ and the difference-path drift rate $R$ enter the second-order planned-time penalty through $\sqrt{\tau_m R/\kappa}$. A valid controller settling bound therefore changes a diagnostic cost with an explicit coefficient. Exact finite-history projections are used when only a few task blocks fit.

\subsection{Relation to previous work}
Experiment choice for hypothesis discrimination has classical sequential formulations due to Chernoff \cite{chernoff} and modern adaptive formulations in active hypothesis testing \cite{naghshvar}. Atkinson and Fedorov develop discriminating regression designs \cite{atkinson}. The predetermined calendar fixes the order of experiment choice and worst healthy-path choice; the sharp constant is specific to that order.

Minimax detection of uncertain signals in Gaussian noise \cite{burnashev} and convex Gaussian testing \cite{gjn} supply the statistical foundation. Likelihood profiling and nuisance rejection are established tools \cite{basseville}. The closest-pair test is an established tool; the new calculations concern the pathwise projection and its optimal dependence on a calendar with transition exclusions.

Robust detection and isolation of abrupt and incipient faults in uncertain nonlinear systems is an established control topic \cite{zhang2002}. Guaranteed active fault detection uses separating inputs for bounded uncertain models \cite{nik98,campbell,scott}; invariant residual sets and support-function duality also lead to active fault-isolation designs \cite{blanchini}. Heirung and Mesbah review this input-design perspective \cite{heirung}. Our noise-free limit is set separation in a two-task experiment. Joint stochastic control and diagnosis \cite{simandl} and the dependence of feedback benefits on the chosen criterion \cite{ashari} further motivate stating the controller and measurement interface separately.

Trend-free run orders with level-change costs \cite{coster}, autoregressive-error run orders \cite{cheng}, and correlated-error sampling designs \cite{bickel1,bickel2} are direct design antecedents. Radiometric timing also balances drift and dead time \cite{schieder,ossenkopf}. Here a transition changes the feasible increment of an unknown mean path, rather than only a counting cost or a prescribed random-error covariance. Independent optimization of each holding window would give a different adversary.

Sensor scheduling for dynamical estimation \cite{vitus,zhao}, observation control and switching costs in quickest change detection \cite{banerjee,lautay,lubenia}, and nonstationary information-growth analysis \cite{liang,hou} address related resource and time constraints. The present result is the sharp second-order deficit under a common rate-limited healthy-path class.

\subsection{Contributions and organization}
The three main results give constant-amplitude rates, sharp growth deficits, and finite-box switching advantages. An admission-design proposition couples a verified settling envelope with the protected statistical score and reduces the tolerance choice to finite calendar comparisons. In the slow-ramp example, the original declared readout standard deviation of $0.01$ N m already gives protected power above $0.999$ at 75 s under symmetric switching, while staying has zero information until about 110 s. Fixed-hold comparisons quantify the extra benefit of the growing-hold rule. A faster-ramp example then shows why a small residual norm alone does not guarantee a useful protected test. Physical calibration and information-interface assumptions are stated in Section~\ref{sec:scope}.
